%! Logarithmic Function Manipulation — Unit 5, Lesson 5.4 — Guided Notes
\documentclass[10pt]{article}
\usepackage{precalculus-article}
\usepackage{precalculus-boxes}
\newcommand{\TallMath}[1]{$\displaystyle #1\rule[-1.4em]{0pt}{3.2em}$}

\begin{document}

\pageheader{Unit 5, Lesson 5.4}{Guided Notes}
\namedateperiod

\vspace{0.1in}

\begin{objectivebox}
\textbf{LO 2.12.A:} Rewrite logarithmic expressions in equivalent forms.

By the end of this lesson I will be able to:
\begin{itemize}[leftmargin=1.5em,itemsep=3pt,topsep=3pt]
  \item Apply the product property to rewrite $\log_b(xy)$ as a sum and connect it to a
        vertical translation of the graph.
  \item Apply the power property to rewrite $\log_b(x^n)$ as a product and connect it to a
        vertical dilation of the graph.
  \item Use the change-of-base formula to evaluate and rewrite logarithms in any base.
  \item Evaluate natural logarithm expressions using $\ln(x) = \log_e(x)$.
\end{itemize}
\end{objectivebox}

\vspace{0.08in}

\begin{vocabbox}
\termblanklong{product property for logarithms}
\termblanklong{power property for logarithms}
\termblanklong{change of base property}
\termblanklong{natural logarithm ($\ln$)}
\termblanklong{vertical translation}
\termblanklong{vertical dilation}
\end{vocabbox}

\vspace{0.08in}

\begin{hookbox}
\textbf{Consider:} We know $\log_2(8) = 3$ and $\log_2(4) = 2$.

\vspace{4pt}
Predict: $\log_2(8 \cdot 4) = \log_2(32) = $ \writeline

\vspace{2pt}
Can you write it as a sum of the two values above?  \writeline
\end{hookbox}

\vspace{0.08in}

\begin{notesbox}{Product Property (EK 2.12.A.1)}
\textbf{Rule:} $\log_b(xy) = \log_b(x) + \log_b(y)$

\vspace{6pt}
\textbf{Graphical connection:} $f(x) = \log_b(kx) = \log_b(k) + \log_b(x)$, so a horizontal
dilation of a logarithmic function is equivalent to a \blank{4cm} of the graph.

\vspace{8pt}
\textit{Practice.} Fill in the blanks.

\vspace{4pt}
$\log_2(8 \cdot 4) = \log_2(\blank{1.5cm}) + \log_2(\blank{1.5cm}) = \blank{1cm} + \blank{1cm} = \blank{1.5cm}$

\vspace{6pt}
$\log_2(32x) = \log_2(\blank{1.5cm}) + \log_2(x) = \blank{1cm} + \log_2(x)$

\vspace{4pt}
(The constant $a = $ \blank{1cm} represents a vertical translation \blank{2cm} by $a$ units.)
\end{notesbox}

\vspace{0.08in}

\begin{notesbox}{Power Property (EK 2.12.A.2)}
\textbf{Rule:} $\log_b(x^n) = n \cdot \log_b(x)$

\vspace{6pt}
\textbf{Graphical connection:} $f(x) = \log_b(x^k) = k\log_b(x)$, so raising the input to a
power is equivalent to a vertical \blank{3.5cm} of the graph by factor $k$.

\vspace{8pt}
\textit{Practice.} Fill in the blanks.

\vspace{4pt}
$\log_3(9^2) = \blank{1cm} \cdot \log_3(9) = \blank{1cm} \cdot \blank{1cm} = \blank{1.5cm}$

\vspace{6pt}
$\log_3(x^4) = \blank{1cm} \cdot \log_3(x)$ \qquad (vertical dilation by factor \blank{1cm})
\end{notesbox}

\vspace{0.08in}

\begin{notesbox}{Change of Base (EK 2.12.A.3)}
\textbf{Rule:} $\log_b(x) = \dfrac{\log_a(x)}{\log_a(b)}$ for any base $a > 0$, $a \ne 1$.

\vspace{6pt}
\textbf{Implication:} Every logarithmic function $\log_b(x)$ is a vertical dilation of $\log_{10}(x)$
by factor $\dfrac{1}{\log_{10}(b)}$. All logarithmic functions are \blank{5cm} of each other.

\vspace{8pt}
\textit{Practice.}

\vspace{4pt}
$\log_5(25) = \dfrac{\log(\blank{1.5cm})}{\log(\blank{1.5cm})} = \dfrac{\blank{1.5cm}}{\blank{1.5cm}} = \blank{1.5cm}$
\end{notesbox}

\vspace{0.08in}

\begin{notesbox}{Natural Logarithm (EK 2.12.A.4)}
The \textbf{natural logarithm} is $\ln(x) = \log_e(x)$, where $e \approx 2.718$.

\vspace{4pt}
It is a logarithmic function with base $e$; all logarithm rules apply to $\ln$.

\vspace{6pt}
Evaluate (no calculator needed):

\vspace{4pt}
$\ln(e) = \log_e(e) = $ \blank{2cm} \qquad $\ln(e^3) = \blank{1.5cm} \cdot \ln(e) = \blank{1.5cm} \cdot \blank{1cm} = \blank{1.5cm}$
\end{notesbox}

\end{document}
